\section{A conductor-forty center trace obstruction}
\label{sec:r671-center-trace}

\begin{theorem}\label{thm:r671-center-trace}
The ring $R8\_2\_671$ has no complex fusion categorification.
\end{theorem}
\begin{proof}
Proposition~\ref{prop:hi6-pair-galois-pu} reduces a hypothetical
categorification to a pseudo-unitary one. Put $q=\sqrt{10}$. Its simple
dimensions and global dimension are
\[
 d=(1,1,3+q,4+q,4+q,3+q,2+q,2+q),\qquad D=120+36q.
\]
The complete ordinary formal codegrees are
$D,D',4,6,6,12,12,12$, where $D'=120-36q$.
The six remaining characters are defined over $\mathbb Q(\sqrt3)$.
All eight source objects are self-dual, and the odd-isotypic pivotal
identities force all their second indicators to be $+1$.
The source induction row is
\[
 [FI(1)]=(8,2,4,6,8,6,6,6).
\]

Write a central forgetful row as $v$, with
$\dim FZ=A+Bq$, where
\[
 A=v_0+v_1+3v_2+4v_3+4v_4+3v_5+2v_6+2v_7,
 \qquad B=\sum_{i=2}^7v_i.
\]
Choose a cyclotomic automorphism $\sigma$ flipping $q$ and fixing
$\sqrt3$. The character-labelled correspondence
\cite[Example~2.9 and Theorem~2.13]{OstrikPiv3} and modular Galois
symmetry give the exact transport
\begin{equation}
 T(v)=v+(0,0,1,1,1,1,1,1)\left(A-\frac{10B}{3}\right).
 \label{eq:r671-transport}
\end{equation}
An actual row has partner $T(v)$ or $-T(v)$, with the sign making it
nonnegative; the entries must be integral. In particular $3\mid B$.
The unit row is paired with $(1,0,1,1,1,1,1,1)$.

Let $J$ be the Galois image of the central unit and put
$w=e_1-e_J$. The multiplicity vector of $I(1)$ is fixed by $S$, so
its Galois signs are positive. Its known twist-one block gives
\[
 \operatorname{proj}_{I(1)}Sw
 =\left(\frac1D-\frac1{D'}\right)w
 =-\frac{w}{\sqrt{40}}.
\]
The torus character of $w$ is the primitive quadratic character
$\chi_{40}$. The equality refinement of the compression bound
\cite[Lemma~4.1]{SchDimensionBlock} implies
$PSw=-w/\sqrt{40}$ and level $40$ for its cyclic representation.
The quadratic Gauss-sum refinement gives negative $\chi_{40}$ on
primitive $40$th twist exponents. These refinements, with their full
local proofs, are supplied in the common compression evidence packet.
Since the $Z$-coordinate of $Sw$ is
$z/D-\sigma(z/D)$, every central simple outside $I(1)$ with
nonrational $z/D$ has twist order dividing $40$ and cannot have twist
one. The last assertion uses the exact full projection $PSw$, whose
right side has no such coordinate.

Let $t$ denote normalized relative trace to $\mathbb Q(q)$.
For a root of unity of arbitrary order, the pair
$(t(\theta),t(\theta^2))$ lies in the convex hull of
\[
 (1,1),\ (-1,1),\ (\tfrac12,-\tfrac12),\
 (-\tfrac12,-\tfrac12),\ (0,-1),\
 (\pm q/8,0),\ (0,\pm q/8).
\]
This is a root-order independent statement: rational traces are
Ramanujan sums; an irrational first trace requires order $40m$, with
$m$ squarefree and coprime to $10$, and has absolute value
$q/(8\varphi(m))$ and second trace zero. An irrational second trace
instead requires order $80m$ and has first trace zero. The displayed
segments contain all these values.
For a high row with nonrational dimension ratio, the allowed pair
list reduces to
\[
 (-1,1),\ (0,-1),\ (-\tfrac14,-\tfrac14),\ (0,0),\
 (\tfrac14,-\tfrac14),\ (0,\tfrac14),\ (-q/8,0).
\]
These are the orders $2,4,5,8,10,20,40$, respectively.
If $v_1>0$, positivity and
$\operatorname{Tr}(\theta^2I(X_1))=D=\dim I(X_1)$ force
$\theta^2=1$. Thus only the two signs remain, and only $-1$ remains
when the ratio is nonrational. Galois partners have equal relative
moments. On the restricted rows one may also choose $\sigma$ with
exponent $59$ modulo $120$; its square is one modulo $40$, so their
actual twists agree.

Here is the complete finite arithmetic needed to finish the proof.
The source tensor determines
\[
 C_{ij}=\sum_{k,l}N_{ki}^{l}N_{kj}^{l}=\sum_Zv_i(Z)v_j(Z).
\]
Each constituent $Z_\chi$ of $I(1)$ has unit multiplicity and
known dimension $D/f_\chi$. Its actual character label imposes
\[
 \frac{\chi_i(v_j)}{f_i}=\frac{\chi_j(v_i)}{f_j}.
\]
After the two forced rows, the six exact bounded domains have sizes
$86,74,74,22,22,22$. Their prescribed sum and the entrywise and
positive semidefinite residual Gram conditions leave exactly six
frames. For every remaining row, integral closure of the modular
formal-codegree relation forces $D/z$ integral; equivalently,
\[
 A^2-10B^2\mid120A-360B,\qquad
 A^2-10B^2\mid36A-120B.
\]
Together with~\eqref{eq:r671-transport} and the residual Gram bound
for the whole partner orbit, the six frames have respectively
$73,85,73,28,36,28$ allowable row orbits.

Combine their Gram columns with the first and second induction trace
identities
\[
 \operatorname{Tr}(\theta I(X_i))=0,\qquad
 \operatorname{Tr}(\theta^2I(X_i))=D\quad(i>0).
\]
After subtracting the known twist-one terms and multiplying the
trace coordinates by eight, these are $56$ integer equations.
The moment domains above produce respectively
$360,398,360,138,168,138$ columns. An actual center would express the
target as a nonnegative real combination of them.

Two fixed integer covectors contradict all six systems. Their
pairings with every allowable column are nonnegative; their target
pairings are $-225800939520$ for the first three frames and $-960$
for the last three. The covectors are given explicitly in
\texttt{r671\_center\_trace\_certificate.json}.
The standalone verifier reconstructs every domain, tests every
rational positive semidefinite condition, and checks every integer
pairing directly from the source tensor. It uses no optimizer or
floating-point arithmetic. Thus every possible center is excluded,
and the initial Galois reduction proves the general complex claim.
\end{proof}
